%% ============================================================================
%% ch1_intro.tex -- Chapter 1: מבוא ורקע   (budget: body pages 1-3, 3 pages)
%% Funnel (booklet fig. 3.1): macro trends -> applied context -> motivation -> roadmap
%% ============================================================================
\chapter{מבוא ורקע}
\label{ch:intro}

\begin{managerbox}
כל החלטת רכישה של מוצר ביוטי או בגד באתר מקוון היא ניחוש של התאמה: האם הגוון יתאים לעור, והאם הגזרה תתאים לגוף. ניחוש שגוי מסתיים בהחזרה, בעלות לוגיסטית ובאובדן אמון. הפרק מראה מדוע הכלים הקיימים, כמו טבלאות מידות, שאלונים ושכבות מציאות רבודה, אינם סוגרים את הפער, ומדוע עכשיו, עם מודלי גוף תלת-ממדיים ומודלי דיפוזיה, אפשר לסגור אותו.
\end{managerbox}

הפרק בנוי כמשפך (איור~\ref{fig:funnel}). הוא פותח במגמות הטכנולוגיות והכלכליות הרחבות, עובר להקשר היישומי של שיתוף הפעולה בין מעבדה אקדמית לקמעונאי גלובלי, מעמיק במוטיבציה ההנדסית והכלכלית, ונחתם במפת הקריאה של הדוח. מבנה זה תואם את מתודולוגיית הדוח המכין של הקורס \cite{segal2026booklet}.

\begin{reportfigure}{פרק המבוא כמשפך: מהמגמה העולמית אל מבנה הדוח}{fig:funnel}
\begin{tikzpicture}[x=1cm,y=1cm]
  \def\lay#1#2#3#4#5{%
    \fill[#5] (-#1,#3) -- (#1,#3) -- (#2,#3-1.05) -- (-#2,#3-1.05) -- cycle;
    \node[font=\small\bfseries,text=white] at (0,#3-0.52) {#4};}
  \lay{5.8}{4.7}{0}{\heb{מגמות: מסחר מקוון, מודלי גוף תלת-ממדיים, דיפוזיה, רגולציה ביומטרית}}{brandprimary}
  \lay{4.7}{3.6}{-1.15}{\heb{ההקשר: מעבדה סטטיסטית וקמעונאי ביוטי ואופנה גלובלי}}{brandprimary!78!brandsecondary}
  \lay{3.6}{2.5}{-2.3}{\heb{המוטיבציה: למה הכלים הקיימים מוצו}}{brandsecondary}
  \lay{2.5}{1.4}{-3.45}{\heb{מבנה הדוח}}{brandaccent!90!black}
  \node[font=\scriptsize,text=brandmuted,anchor=west] at (4.9,-0.5) {\heb{רחב}};
  \node[font=\scriptsize,text=brandmuted,anchor=west] at (2.6,-4.0) {\heb{ממוקד}};
  \draw[-Stealth,brandmuted,thick] (5.3,-0.9) -- (3.1,-3.9);
\end{tikzpicture}
\end{reportfigure}

מהאיור עולה שכל שכבה צרה מקודמתה, וכל טענה בשכבה התחתונה צריכה להישען על שכבה רחבה יותר שמעליה. לכן הפרק אינו נפתח בקביעה כללית על כוחה של הבינה המלאכותית, אלא בשלושה תהליכים מדידים שמתרחשים בו-זמנית.

\section{מגמות-על טכנולוגיות, כלכליות ורגולטוריות}
\label{sec:trends}

המגמה הראשונה היא מעבר של קטגוריות שבעבר נקנו כמעט רק בחנות, כמו מייקאפ, צבע שיער ובגדים מחויטים, אל ערוצים דיגיטליים. בקטגוריות אלה ההחלטה תלויה בתכונות גוף שאינן נראות במסך: גוון עור, תת-גוון (\en{Undertone}), צבע עיניים ופרופורציות גוף. לפי אומדני הצוות, המבוססים על סדרי גודל שמקובלים בתעשייה וייבדקו מול דוחות אנליסטים בשלב המימוש, שיעור ההחזרות בהלבשה מקוונת עומד על \num{30\%--40\%}, ורובו נובע מאי-התאמה של מידה וגזרה. סקירות עדכניות של מערכות המלצה לאופנה מצביעות על אותו פער מהצד האלגוריתמי: רוב המערכות לומדות מהיסטוריית קליקים ורכישות, ולא ממאפיינים פיזיים מכוילים של הלקוחה \cite{deldjoo2023fashionrec}.

המגמה השנייה טכנולוגית. בתוך פחות מחמש שנים התבגרו שלוש משפחות כלים: מודלי גוף פרמטריים כמו \en{SMPL-X}, שמשחזרים גוף, ידיים ופנים מתמונה יחידה \cite{pavlakos2019smplx,goel2023hmr2}; מודלי דיפוזיה סמויה (\en{Latent Diffusion}) שמייצרים תמונות פוטו-ריאליסטיות בעלות חישובית סבירה \cite{rombach2022ldm}; ומודלי שיכון רב-מודאליים כמו \en{CLIP}, שממפים תמונה וטקסט לאותו מרחב וקטורי \cite{radford2021clip}. השילוב של שלושתם מאפשר לראשונה לבנות "תאום דיגיטלי" של המשתמשת, שעליו אפשר להלביש, לאפר ולהמליץ.

המגמה השלישית רגולטורית. תמונת פנים, צבע עיניים ומידות גוף הם נתונים ביומטריים או נתונים שמאפשרים זיהוי, ולכן חלות עליהם דרישות מחמירות של צמצום מידע והסכמה מפורשת לפי תקנת \en{GDPR} \cite{gdpr2016}. מגמה זו אינה רק אילוץ: היא דוחפת את התכן לכיוון של עיבוד בקצה (\en{Edge}), שבו התמונה הגולמית אינה עוזבת את המכשיר. חוק הבינה המלאכותית האירופי (\en{AI Act}) מוסיף שכבה נוספת: הוא מסווג מערכות לפי רמת סיכון, ומטיל חובות שקיפות ותיעוד על מערכות שמעבדות מאפיינים ביומטריים \cite{euaiact2024}. מערכת המלצות קוסמטית אינה מערכת זיהוי, אך תכן שמפריד בין כיול לבין זיהוי, ושומר רק וקטורי צבע ומידות במקום תמונות, מקטין מראש את החשיפה הרגולטורית.

לצד הרגולציה, כוח החישוב במכשיר הנייד גדל. מכשירי ביניים כוללים כיום יחידת עיבוד עצבי (\en{NPU}) שמריצה מודלי פילוח קלים ואיתור נקודות ציון בזמן אמת, כפי שהודגם במערכות כמו \en{MediaPipe Hands} \cite{zhang2020mediapipe}. המשמעות ההנדסית היא שאפשר לפצל את החישוב: מה שרגיש ומהיר נשאר במכשיר, ומה שכבד ונדיר עובר לענן. החלטת הפיצול היא אחת ההחלטות המרכזיות שייבחנו בפרק~\ref{ch:design}.

\section{ההקשר היישומי: מעבדה סטטיסטית וקמעונאי גלובלי}
\label{sec:context}

הפרויקט מתוכנן כשיתוף פעולה בין צוות מהחוג לסטטיסטיקה באוניברסיטת חיפה לבין שותף תעשייתי מדומה, שמייצג מעבדת חדשנות של קבוצת קוסמטיקה וקמעונאית אופנה גלובלית. השותף מוצג בדוח כמודל עבודה בלבד, ואין בו טענה להסכם קיים. תפקיד השותף הוא לספק קטלוג מוצרים עם מאפייני צבע ומידות לכל יחידת מלאי (\en{SKU}), ממשק שותפים (\en{Affiliate}) למדידת המרות, ומדגם פיילוט של משתמשות. תפקיד המעבדה הוא לספק את התשתית הסטטיסטית: אמידת אי-ודאות, תכנון ניסויים, ומדדי אימות שאפשר להגן עליהם מול ועדה.

טבלה~\ref{tab:roles} מפרטת את חלוקת התרומות. ממנה עולה שהשותף התעשייתי תורם נתונים וערוץ הפצה, והמעבדה תורמת את שכבת המדידה והאימות. לכן הפרויקט אינו תלוי בשותף לצורך הוכחת ההיתכנות המחקרית: נתוני הכיול נאספים בניסוי מעבדה, והקטלוג יכול להיות מוחלף בקטלוג ציבורי בתרחיש מגירה.

\begin{reporttable}{חלוקת התרומות בין המעבדה האקדמית לשותף התעשייתי (המדומה)}{tab:roles}{P{2.6cm} R R}
\headrow \hc{תחום} & \hc{המעבדה (החוג לסטטיסטיקה)} & \hc{השותף התעשייתי} \\ \midrule
נתונים & ניסוי כיול מבוקר, \num{120} משתתפות & קטלוג \en{SKU} עם ערכי צבע ומידות \\
שיטות & אמידה בייסיאנית, תכנון ניסויים, ניתוח עוצמה & ממשק שותפים ומדידת המרות \\
תשתית & סביבת פיתוח, מאגר \en{Git}, מנוע סימולציה & גישת \en{API} לקטלוג ולמלאי \\
אימות & מבחני השערות, רווחי סמך, ניתוח רגישות & פאנל סטייליסטיות ופיילוט משתמשות \\
\end{reporttable}

\begin{examplebox}[יישום במערכת \en{AuraTwin AI}]
משתמשת מצלמת סלפי ומדווחת שצבע עיניה חום. המערכת מזהה את לובן העין ואת הקשתית, משתמשת בהם כעוגני צבע, ומתקנת את צבע העור לפני שהיא בוחרת גוון מייקאפ. במקביל היא אומדת היקפי חזה, מותן ואגן ממודל גוף, מצליבה עם טרנד עונתי שזוהה ברשתות החברתיות, ומציעה שלושה פריטים עם קישורי שותפים. הדמיה פוטו-ריאליסטית של הפריט על גופה נפתחת לאחר צפייה במודעה מתגמלת.
\end{examplebox}

\section{המוטיבציה ההנדסית והכלכלית: מדוע הכלים הקיימים מוצו}
\label{sec:motivation}

שלוש משפחות פתרונות שולטות היום בשוק, ולכל אחת מגבלה מבנית שאינה ניתנת לתיקון בשיפור הדרגתי:

\textbf{טבלאות מידות ושאלונים סטטיים.} הלקוחה מתבקשת למדוד את עצמה או לבחור "סוג גוף" מתוך רשימה. השיטה זולה, אך היא מעבירה את שגיאת המדידה אל הלקוחה, ואינה מטפלת כלל בצבע. שאלון שמבקש לבחור "גוון עור" מתוך שש דוגמאות על מסך לא מכויל משלב שתי שגיאות: של המסך ושל הזיכרון.

\textbf{שכבות מציאות רבודה (\en{AR}).} מערכות אלה מציירות שפתון או צבע שיער על פנים בזמן אמת, ונותנות חוויה מרשימה. אולם הן מציירות את הגוון המבוקש מעל תמונה שכבר עוותה בידי איזון הלבן של המצלמה. כפי שמראה ספרות תיקון איזון הלבן, כשל באיזון הלבן יוצר שגיאה גלובלית ולא-לינארית שאינה מתוקנת בפשטות \cite{afifi2019wb}. סימולציה מקדימה של הצוות (פרק~\ref{ch:problem}) מראה שבתמונת סלפי, שבה העור תופס את רוב השטח, הנחת "העולם האפור" (\en{Gray World}) מחמירה את השגיאה במקום לתקן אותה.

\textbf{הדמיה גנרטיבית בענן.} מודלי הדמיית לבוש מבוססי דיפוזיה הגיעו לאיכות פוטו-ריאליסטית \cite{zhu2023tryondiffusion,choi2024idmvton}, אך כל תמונה דורשת עשרות צעדי ניקוי רעש על מעבד גרפי (\en{GPU}). כשהשירות ניתן בחינם ובלי הגבלה, עלות החישוב גדלה לינארית עם השימוש, בעוד ההכנסה אינה גדלה איתה. זו בעיה כלכלית, לא אלגוריתמית, והיא מחייבת שהמודל העסקי יהיה חלק מהארכיטקטורה עצמה.

מנקודת מבט של מנהל, שלוש המגבלות מתלכדות לחשבון אחד. כל החזרה של מוצר נושאת עלות של משלוח חוזר, בדיקה ואריזה מחדש, ומוצרי קוסמטיקה שנפתחו לרוב אינם חוזרים למלאי כלל. לכן שיפור קטן בהסתברות ההתאמה שווה יותר משיפור גדול בחוויית הממשק. גם בצד ההכנסות התמונה ברורה: המלצה שמגיעה עם קישור שותפים מניבה עמלה רק אם היא מדויקת מספיק כדי להוביל לרכישה שאינה מוחזרת. דיוק המדידה הוא אפוא משתנה כלכלי, ולא רק מדד איכות טכני.

\begin{keymessagebox}
הפער אינו בהיעדר מודל טוב יותר לאחד הרכיבים, אלא בהיעדר מערכת שמחברת כיול צבע ביומטרי, מידות גוף מכוילות, אחזור טרנדים והדמיה גנרטיבית, תחת מודל כלכלי שמממן את החישוב.
\end{keymessagebox}

\section{מבנה הדוח}
\label{sec:structure}

פרק~\ref{ch:problem} מגדיר את הבעיה במספרים, מנסח את שאלת המחקר ואת יעדי התכן, ומציג את מפרט הדרישות ההנדסי. פרק~\ref{ch:lit} סוקר את הספרות בארבעה תחומי ידע, בוחר את מאמר הדגל ומנסח את פער הידע. פרק~\ref{ch:market} מנתח את היקף השוק ואת המתחרים מנקודת מבט טכנית. פרק~\ref{ch:design} משווה שלוש ארכיטקטורות במטריצת החלטה משוקללת, ומפרט את הארכיטקטורה שנבחרה ואת חדשנותה. פרק~\ref{ch:plan} מציג את תוכנית העבודה ל-\num{14} שבועות, פרק~\ref{ch:risk} את ניתוח הסיכונים, ופרק~\ref{ch:summary} מסכם את מוכנות הפרויקט לשלב המימוש. הדוח הופק בתבנית לטך דו-לשונית בשתי שכבות, שנבנתה בשיטת "מתארים במילים, המודל כותב, הסטודנטיות חותמות" \cite{segal2026l03}. התבנית, הפרומפטים וסוכני הבדיקה זמינים במאגר הפרויקט, כך שכל עמוד בדוח ניתן לשחזור.
